We consider the supervised binary classification setting studied in~\cite{lawless_interpretable_2023}.
%We refer to this source as the ``original paper''.
What follows is a recap that closely follows this paper.
Given is a training set of $n$ points, denoted by $\I$, from an underlying (unknown) distribution.
Each point has $m$ features from set $\J$, plus a label \[(\X_i, y_i), \quad i \in \I = \{1, \dots, n\}.\]
In this setting, feature values $\X_i \in \{0,1\}^\J$ and labels
$y_i \in \{-1,1\}$ for all $i \in \I$.
A point $i$ is said to possess feature $j$ if $X_{ij} = 1$, where $X_{ij}$ is the $j$-th element of vector $\X_{i}$.
A point $i$ is said to be positive if $y_i = 1$ and negative otherwise.

Theoretically, binary features are not restrictive.
For example, categorical features can be expressed as binary features via one-hot encoding.
Real-valued features can be discretized into a number of binary features via thresholds,
i.e., binary features reflecting $X_{ij} \leq v$ or $X_{ij} \geq v$ may be introduced for various $v$.
The number of features can grow rapidly with this approach.

The objective of this model is to learn a ruleset as described earlier.
Mathematically, define the set of possible clauses as $\cL = \J$.
A rule is a subset of $\cL$.
The candidate rule set is given by  $\K = 2^\L$, i.e.\ the power set of $\L$.
Let $\K_i$ be the set of rules satisfied by point $i$, so
\begin{equation}
    \K_i = \left\{k \in \K : \underset{j \in k}{\bigwedge} X_{ij} \right\}, \quad i \in \I. \label{eq:boolean_K_i}
\end{equation}
In other words, $\K_i$ is the set of rules that cover point $i$.
Observe that $\emptyset \in \K$ and that the conjunction in \eqref{eq:boolean_K_i}
is always true for $k=\emptyset$.
So $\emptyset \in \K_i$ for any $i \in \I$.
Therefore, the ``empty rule'', $\emptyset$, covers all points under this definition.



\begin{definition}[DNF Classifier, \cite{lawless_interpretable_2023}]\label{def:dnf_classifier}
    Given a ruleset $S \subseteq \K$, the DNF classifier is defined as follows:
    \[\hat{y}_S(\textbf{X}_i) ~=~
    \begin{cases}
        1, & \text{if } \abs{\K_i \cap S} > 0, \\
        0, &\text{else.}
    \end{cases}\]
\end{definition}

There exist $2^m$ possible rules, i.e. $|\K| = 2^m$.
Therefore, as $m$ grows, enumerating all possible rules quickly becomes intractable.
In practice, one may wish to restrict $S$, for example by limiting the number of clauses to $D$.
This restricts the number of possible rules to \[\sum_{d=0}^D \binom{m}{d}.\]


The aim is to minimize the \emph{\zo classification error}, which counts the number of datapoints that are misclassified by a rule set $S$.
Hence, it is the sum of the \emph{\zo loss} of each datapoint.

\begin{definition}[\zo loss and \zo classification error]\label{def:01_loss}
    The \emph{\zo loss} of point $i$ subject to ruleset $S$ is
    \[\ell_S(\bm{X}_i) = \Ind\left(\hat{y}_S(\bm{X}_i) \neq y_i\right) =
    \begin{cases}
        \Ind(\abs{\K_i \cap S }= 0), & \text{if } y_i = 1,\\
        \Ind(\abs{\K_i \cap S} > 0), &\text{otherwise},
    \end{cases}\]
    where $\Ind(\mathcal{A})$ is the indicator function (i.e., $\Ind(\mathcal{A}) = 1$ if $\mathcal{A}$ is true and 0 otherwise).
The \emph{\zo classification error} of ruleset $S$ is
    \[\ell_S(\bm{X}) =
    \sum_{i \in \I} \ell_S(\bm{X}_i). \]
\end{definition}
Closely related is the accuracy, which is the fraction of correctly classified points.
\begin{definition}[Accuracy]
    The \emph{accuracy} of a ruleset $S$ is
    \[
        \text{Accuracy}(S) = \frac{1}{n}\left|\left\{i \in \I : \hat{y}_S(\bm{X}_i) = y_i\right\}\right|.
    \]
\end{definition}
Minimizing the \zo classification error is equivalent to maximizing the classification accuracy, as illustrated in
the introduction to this chapter.


We now formulate the problem of finding a ruleset with minimum \zo classification error as an Integer Program (IP).
Partition the points into two sets, the positive class and the negative class:
\[\cP = \{i \in \I : y_i = 1\} \text{ and } \cN = \{i \in \I : y_i \neq 1\}.\]
Let $w_k \in \{0,1\}$ be a variable indicating if rule $k \in \K$ is used in the rule set.
Let $\zeta_i\in\{0,1\}$  be a variable indicating if point $i\in \pun$ is misclassified.
The IP is given by


%\begin{alignat}{3}
%\textbf{min} && \quad \sum_{i\in \cal{P}} \zeta_i +\sum_{i\in \cal{N}} \zeta_i ~~ \label{obj:01}  \\
%	\textbf{s.t.}  && \quad
%	 \zeta_i + \sum_{k\in \K_i} w_k \geq & 1 , \quad && \forall i \in \cal{P}, \label{const:accP01} \\
%	 && w_k  \leq &   \zeta_i , \quad  && \forall i \in {\cal N}, k \in \K_i,  \label{const:accZ01} \\[.1cm]
%	%	& w_k,  ~\zeta_i  \in\{0,1\}, \quad \forall {k\in \cal{K}}, i \in {\cal P} \cup {\cal N}. ~~ \label{const:accBin01} \\
%	&& w\in \{0,1\}^{\abs{\K}},~ & \zeta\in \{0,1\}^{\abs{\cP \cup {\cal N}}} \label{const:accBin01}
%\end{alignat}

\begin{subequations}
    \label{eq:BM_orig}
    \begin{align}
        \mathbf{\min} \quad & \sum_{i \in \cP} \zeta_i + \sum_{i \in \cN} \zeta_i \label{eq:BM_orig-obj} \\
        \st \quad & \zeta_i + \sum_{k \in \K_i} w_k \geq 1, && \forall i \in \cP, \label{eq:BM_orig-misP} \\
        & w_k \leq \zeta_i, && \forall i \in \cN,\; k \in \K_i, \label{eq:BM_orig-misN} \\
        & w \in \{0, 1\}^{\K}, \quad \zeta \in \{0, 1\}^{\cP \cup \cN}. \quad && \label{eq:BM_orig-vars}
    \end{align}
\end{subequations}
The objective of this IP is exactly the \zo classification error.
To see that this IP is correct, consider first \eqref{eq:BM_orig-misP}.
Let $i \in \P$.
If $w_k = 1$ for some $k \in \K_i,$ then $i$ is covered by the ruleset.
Indeed, then $\zeta_i=0$ is valid and optimal.
If not, then $i$ is not covered and $\zeta_i = 1$ to satisfy the constraint.
In either case, $\zeta_i$ correctly represents whether $i$ is misclassified.
Now consider \eqref{eq:BM_orig-misN}.
If $w_k = 1$ for some $k \in \K_i$, then $i$ is covered and \eqref{eq:BM_orig-misN} implies $\zeta_k = 1$.
Otherwise, $\zeta_i =0$ is valid and optimal.
In either case, $\zeta_i$ correctly represents whether $i$ is misclassified.


This formulation has one integer variable per point and multiple constraints per negative datapoint.
For a large number of points, the IP becomes large.
In many practical classification tasks, the number of positive points is small compared to the total.
To limit the size of the IP, the Hamming loss may be considered.
In the Hamming loss, we do not count the number of negative points that are misclassified, but we count
the total number of times a negative point is misclassified by a rule.

\begin{definition}[Hamming loss \cite{lawless_interpretable_2023}] The \emph{Hamming loss} of point $i$ under ruleset $S$ is given by
\[\ell^{h}_S(\bm{X}_i) =
\begin{cases}
    \Ind(\abs{\K_i \cap S }= 0), & \text{if } y_i = 1,\\
    \abs{\K_i \cap S}, &\text{otherwise.}
\end{cases}\]
\end{definition}

It has been illustrated that Hamming loss is a good alternative to \zo loss in practical settings~\cite{lawless_interpretable_2023}.

Furthermore, it is useful in practice to restrict the complexity of a ruleset.
Define the complexity of a rule $k \in \K$ as
\begin{equation}
    c_k = 1 + |k|, \label{eq:boolean_ck}
\end{equation}
so one plus the number of features in the rule.
The complexity of a ruleset $S \subseteq \K$ is given by $\sum_{k\in S} c_k,$ which we bound by $C$.

The new IP that uses Hamming loss becomes

\begin{subequations}
    \label{eq:BM}
    \begin{align}
        \mathbf{\min} \quad & \sum_{i \in \cP} \zeta_i + \sum_{i \in \cN} \sum_{k \in \K_i} w_k \label{eq:BM-obj} \\
        \mathbf{\text{s.t.}} \quad & \zeta_i + \sum_{k \in \K_i} w_k \geq 1, & \forall i \in \cP, \label{eq:BM-misP} \\
        & \sum_{k \in \K} c_k w_k \leq C, & \label{eq:BM-complex} \\
        & w \in \{0, 1\}^{\K}, \zeta \in \{0, 1\}^{\cP}. \label{eq:BM-binary2}
    \end{align}
\end{subequations}
For a positive point $i$, the justification given for \eqref{eq:BM_orig} still holds.
For a negative point $i$, there is no $\zeta_i$, but instead the number of rules wrongly covering it are counted
directly in the objective.

The number of possible rules $|\K|$ grows quickly as $m$ grows.
Therefore, column generation is used to solve the LP relaxation of this problem.
A simple approach for obtaining an integer solution is employing a price-then-branch heuristic.
In this technique, the relaxed master problem (RMP) is first solved to optimality.
Next, the restricted integer master problem, with only the previously generated columns, is solved to optimality using a standard IP solver.
While it is possible to solve the problem exactly via a branch-and-price approach,
this is harder to implement, it is slower,
and it is unclear whether it improves the quality of the output significantly.
In some cases, the price-then-branch heuristic can reveal that an exact solution is
only marginally better in terms of objective.